%!TEX program = uplatex
\RequirePackage{plautopatch}
\documentclass[uplatex,dvipdfmx]{jlreq}
\usepackage{amsmath,amssymb}

\pagestyle{empty}

\begin{document}

%======================================================================
% 表紙
%======================================================================
\begin{center}
  \vspace*{20pt}
  {\Huge\bfseries ハミルトン力学系}

  \vspace{30pt}

  {\Large\bfseries 天体力学・シンプレクティック幾何学・可積分系}

  \vspace{50pt}

  {\large 前文}

  \vspace{10pt}

  \parbox{0.85\textwidth}{
    \large\bfseries
    ケプラーからポアンカレに至る天体力学の流れを辿りつつ、
    KAM 理論・シンプレクティック・トポロジー・特異点解析という
    現代的視点からハミルトン力学系を多角的に論じる。
  }
\end{center}

\vfill

\newpage

%======================================================================
% 講演１：伊藤秀一
%======================================================================
\noindent
{\large\bfseries ハミルトン力学系への招待}

\medskip
\hspace*{\fill}{\large\bfseries 伊藤 秀一（東工大・理）}

\bigskip\bigskip

天体力学（celestial mechanics）はニュートンの運動法則によって記述される
天体の運動を研究対象とする学問である。
それはケプラー以来の古典力学の華であるとともに、
とくに19世紀はじめまでは数学と一体になって発展し
（たとえばガウスはゲッティンゲンの天文台長でもあった）、
その後もポアンカレにいたるまで、
数学に対して最も大きな影響を及ぼした分野の一つであろう。
そして今回取りあげる、力学系、シンプレクティック幾何学、可積分系といった現代の研究分野も、
もとを辿れば古典力学とくに天体力学の問題が関っていることが多い。

ニュートンは「プリンキピア」（1687年出版）の中で、
ケプラーの法則から惑星が太陽からの距離の２乗に反比例する力を受けることを証明し、
その後ベルヌーイやオイラーらによって微積分法が発展するにつれて、
力学の取り扱いは微分方程式を用いた「解析的」なものへと変貌していく。
とくにラグランジュは、ニュートンの運動方程式を、
配位空間（質点系の位置を表す空間）の座標の取り方によらない、
今日オイラー・ラグランジュ方程式と呼ばれる方程式として書き表し、
運動方程式の取り扱いに進歩をもたらした。
ラグランジュの著作「解析力学」が出版されたのはプリンキピアから約100年後の1788年である。
さらに19世紀前半には、ハミルトンやポアソン、ヤコビらによって
力学の理論形式は整備され、今日「解析力学」と呼ばれる分野の基礎ができたといえる。
ハミルトンは幾何光学に関連した研究からハミルトン--ヤコビ方程式に到達し、
一方ヤコビはそれを整備し、正準変換や母関数の概念を導入するとともに、
ハミルトン--ヤコビ方程式を用いたケプラー問題や楕円面上の測地線の運動の求積などを行った。
その研究は楕円関数の理論の発展などとも相まって、「古典解析」の大きな流れを形成したといえよう。
とくに今日ハミルトン系と呼ばれる「ハミルトンの正準方程式」（これはヤコビの命名である）は、
相空間（配位空間×運動量の空間）の上の座標変換（正準変換）で不変であり、
シンプレクティック幾何の萌芽もこの頃には現れたといえる。

このような力学理論の進歩にもかかわらず、
天体力学の難問「３体問題」は未解決のまま残され、
その進展には19世紀末のポアンカレの研究まで待たなければならなかった。
ポアンカレは、３体問題が求積法では解ける見込みがないこと、
さらには今日「カオス」と呼ばれる現象の原型（ホモクリニック軌道）を発見し、
それらの研究を含む著作「天体力学における新しい方法」（全３巻）は
1892年から1899年にかけて出版されている。
ポアンカレは、それまでの求積法に代表されるような「解析的」手法に比べて、
周期解の存在と Twist 写像の不動点定理を結び付けるといった「幾何的な」手法を導入し
微分方程式の解の定性的研究を生んだが、
その研究を育んだのは制限３体問題という天体力学における典型的問題だったのである。
その研究は力学系理論を創始するとともに、その後の発展に決定的な影響を与えている。

今回の私の話では、今述べたようなケプラーから始まりポアンカレに至る
天体力学あるいは一般の古典力学をめぐる数学の研究の流れを第１部で辿り、
それらを現代の視点から眺めてみたい。
そして第２部では、摂動論とくにポアンカレが力学の基本問題と呼んだ
「可積分系の摂動論」（制限３体問題はその典型例である）について、
コルモゴロフ--アーノルド--モーザー（KAM）理論の概要と
それに付随した様々な問題についてお話したい。

\bigskip

\noindent{\large\bfseries 参考文献}

\begin{enumerate}
  \setlength{\itemsep}{2pt}
  \item F.~Diacu and P.~Holmes,
        \textit{Celestial Encounters --- The Origins of Chaos and Stability},
        Princeton Univ.\ Press, 1996.
  \item J.~Moser,
        \textit{Stable and Random Motions in Dynamical Systems},
        Annals of Math.\ Studies \textbf{77}, Princeton Univ.\ Press, 1973.
  \item 齋藤利弥，解析力学，至文堂，1964．
  \item 伊藤秀一，常微分方程式と解析力学，共立出版，1998．
\end{enumerate}

\newpage

%======================================================================
% 講演２：小野薫
%======================================================================
\noindent
{\large\bfseries ポアンカレ--バーコフの不動点定理からアーノルド予想へ}

\medskip
\hspace*{\fill}{\large\bfseries 小野 薫（北大・理）}

\bigskip\bigskip

天体力学の問題から派生した Poincaré--Birkhoff の不動点定理と
それが一つの出発点となった symplectic topology、
特に Hamilton 系の周期解の存在についていくつかの話題を紹介する予定です。

Poincaré--Birkhoff の不動点定理とは、円環領域の面積・向きを保つ同相写像が
twist 条件と呼ばれるものを満たすとき、少なくとも二つの不動点を持つというものです。
恒等写像に近い場合は、generating function と呼ばれる関数の臨界点の話に帰着されます。
こうした考え方はさらに一般化できて、
余接束の symplectic 幾何学の基本的道具となっています。
また、これは Hamilton 系の action functional と呼ばれるものの有限次元近似を与えています。
一般の空間で考える時には、こうした近似をせずに無限次元でものを考えることになります。
今回の話では、あまり一般の場合は扱わない積りです。
余接束あるいは symplectic vector space の中でも面白い話があることを紹介したいと考えています。

\newpage

%======================================================================
% 講演３：吉田春夫
%======================================================================
\noindent
{\large\bfseries コワレフスカヤの夢——積分可能性と解の一価性の接点}

\medskip
\hspace*{\fill}{\large\bfseries 吉田 春夫（国立天文台）}

\medskip
\hspace*{\fill}〒181--8588　東京都三鷹市大沢 2-21-1

\medskip
\hspace*{\fill}\texttt{yoshida@gauss.mtk.nao.ac.jp}

\bigskip\bigskip

自由度 $n$ の Hamilton 力学系は $n$ 個の包合的な第一積分を許容するとき
積分可能（可積分）と呼ばれる。
積分可能系の解は一般に準周期的となりカオスは出現しえない。
そこで具体的に与えられた系が積分可能か否かを判定することは
Hamilton 力学系の最も基本的な問題の一つであるはずであるが、全く未解決の問題である。
もちろん最終的な判定条件そのもの自身、存在が保証されているわけではない。
しかるに近年この方向でめざましい発展が見られた。
そこでこの判定条件の候補として期待されているのが特異点解析
あるいは Painlevé 解析と呼ばれているものであり、
前世紀末1888年のコワレフスカヤの研究に端を発する。

積分可能系としてのコワレフスカヤのコマの発見は、
Hamilton 系の積分可能性と解の複素時間平面における一価性という解析的な性質との
隠れた関係の存在を示唆した。
非線形の常微分方程式は一般に解の複素時間平面に特異点を持つ。
多くの具体的な例題はこの特異点の性質と系の積分可能性との間にある隠された関係の存在を示唆する。
定理という形での正当化が十分なされないまま、
この特異点解析の「特異点は極のみ」という要請は
既存の多くの積分可能系を特徴づけ、また新たな積分可能系を発見するのに役立ってきた。

本講演の前半ではコワレフスカヤのコマの発見に至る歴史的事実から始め
特異点解析が有効に働く具体例を列挙する。
後半ではこの経験的な特異点解析を正当化する試みをいくつか紹介することにしたい。
最新のものでは線形常微分方程式に対する Galois 理論と位置づけられる
Picard--Vessiot 理論の非線形 Hamilton 系版（Morales--Ramis, 1997）に
Gauss の超幾何方程式の求積による可解性条件を与えた
木村俊房の定理（1969）を適用して得られた結果がある。

\bigskip

\noindent{\large\bfseries 参考文献}

\begin{enumerate}
  \setlength{\itemsep}{2pt}
  \item 大貫義郎・吉田春夫，岩波講座現代の物理学１・力学，岩波書店（1994）．
  \item 吉田春夫，コワレフスカヤのコマ，数理科学，1981年１月号．
\end{enumerate}

\end{document}
